\documentclass[]{article}
\usepackage{lmodern}
\usepackage{amssymb,amsmath}
\usepackage{ifxetex,ifluatex}
\usepackage{fixltx2e} % provides \textsubscript
\ifnum 0\ifxetex 1\fi\ifluatex 1\fi=0 % if pdftex
  \usepackage[T1]{fontenc}
  \usepackage[utf8]{inputenc}
\else % if luatex or xelatex
  \ifxetex
    \usepackage{mathspec}
    \usepackage{xltxtra,xunicode}
  \else
    \usepackage{fontspec}
  \fi
  \defaultfontfeatures{Mapping=tex-text,Scale=MatchLowercase}
  \newcommand{\euro}{€}
\fi
% use upquote if available, for straight quotes in verbatim environments
\IfFileExists{upquote.sty}{\usepackage{upquote}}{}
% use microtype if available
\IfFileExists{microtype.sty}{\usepackage{microtype}}{}
\usepackage[margin=1in]{geometry}
\ifxetex
  \usepackage[setpagesize=false, % page size defined by xetex
              unicode=false, % unicode breaks when used with xetex
              xetex]{hyperref}
\else
  \usepackage[unicode=true]{hyperref}
\fi
\hypersetup{breaklinks=true,
            bookmarks=true,
            pdfauthor={Justin Le},
            pdftitle={Effective LLM-Assisted Haskell: Understanding Constraint-Evading Behavior Draft},
            colorlinks=true,
            citecolor=blue,
            urlcolor=blue,
            linkcolor=magenta,
            pdfborder={0 0 0}}
\urlstyle{same}  % don't use monospace font for urls
% Make links footnotes instead of hotlinks:
\renewcommand{\href}[2]{#2\footnote{\url{#1}}}
\setlength{\parindent}{0pt}
\setlength{\parskip}{6pt plus 2pt minus 1pt}
\setlength{\emergencystretch}{3em}  % prevent overfull lines
\setcounter{secnumdepth}{0}

\title{Effective LLM-Assisted Haskell: Understanding Constraint-Evading Behavior
Draft}
\author{Justin Le}

\begin{document}
\maketitle

\emph{Originally posted on
\textbf{\href{https://blog.jle.im/entry/effective-llm-assisted-haskell-understanding-constraint-evading-behavior-draft.html}{in
Code}}.}

\documentclass[]{}
\usepackage{lmodern}
\usepackage{amssymb,amsmath}
\usepackage{ifxetex,ifluatex}
\usepackage{fixltx2e} % provides \textsubscript
\ifnum 0\ifxetex 1\fi\ifluatex 1\fi=0 % if pdftex
  \usepackage[T1]{fontenc}
  \usepackage[utf8]{inputenc}
\else % if luatex or xelatex
  \ifxetex
    \usepackage{mathspec}
    \usepackage{xltxtra,xunicode}
  \else
    \usepackage{fontspec}
  \fi
  \defaultfontfeatures{Mapping=tex-text,Scale=MatchLowercase}
  \newcommand{\euro}{€}
\fi
% use upquote if available, for straight quotes in verbatim environments
\IfFileExists{upquote.sty}{\usepackage{upquote}}{}
% use microtype if available
\IfFileExists{microtype.sty}{\usepackage{microtype}}{}
\usepackage[margin=1in]{geometry}
\ifxetex
  \usepackage[setpagesize=false, % page size defined by xetex
              unicode=false, % unicode breaks when used with xetex
              xetex]{hyperref}
\else
  \usepackage[unicode=true]{hyperref}
\fi
\hypersetup{breaklinks=true,
            bookmarks=true,
            pdfauthor={},
            pdftitle={},
            colorlinks=true,
            citecolor=blue,
            urlcolor=blue,
            linkcolor=magenta,
            pdfborder={0 0 0}}
\urlstyle{same}  % don't use monospace font for urls
% Make links footnotes instead of hotlinks:
\renewcommand{\href}[2]{#2\footnote{\url{#1}}}
\setlength{\parindent}{0pt}
\setlength{\parskip}{6pt plus 2pt minus 1pt}
\setlength{\emergencystretch}{3em}  % prevent overfull lines
\setcounter{secnumdepth}{0}


\begin{document}

Sooo yes it's true, I've been integrating LLMs and agentic coding tools in my
Haskell coding since the beginning of this year. I remember last year barely
being able to prompt my heavily GADT-ized [fancy
Haskell](https://blog.jle.im/entry/extreme-haskell-typed-expression-edsls-1.html)
Haskell into the web UI for ChatGPT/Gemini, progressing to using Open AI's
web-based codex host, to biting the bullet and finally permitting a sentient
Eldritch being to persistently live and experience existence and qualia on my
personal machines and ship-of-theseus dev infrastructure I've been operating
continuously since freshman year of uni.

I've been using it for a lot of my projects, both personal and professional.
Along the way I've collected some habits and patterns for using them more and
more effectively. I've dropped notes [on various
articles](https://blog.jle.im/entries/tagged/agentic.html) about small things,
but hesitated on aggregating them into an actual blog post...mostly because I'm
"new" at this (six months?), and also because the nature of interacting with
LLMs and their limitations seems to be re-written every few months.

But, after some recent community rumblings and fragmentation regarding
high-stakes existential tensions between agentic LLM coding and the future
direction of Haskell, I realized I could probably lend a constructive hand by
describing how I, personally, a Haskeller of 15+ years, (7+ professional)
personally make the most out of them.

Note: none of the views espoused in this post are endorsed by my employer. They
are my own and will most likely change drastically as I learn to use these tools
better over time, and, as presented, are uniquely formed around my background
and professional skill set and most likely not *directly* applicable to any
arbitrary human. However, I do feel there are people who may appreciate my
personal experience report.

This post will focus on what I call "constraint-evading behavior" in LLMs as it
relates to effective Haskell.

## Ground Rules

First let me clarify at what level I integrate LLM's at.

For any projects of importance:

-   Everything I commit I treat with as much standards line-by-line as any code
    I write.
-   All design decisions and concepts and code structures are written in the
    style and care as if I had written it by hand.
-   No line of code is "less scrutinized" or "abstracted away" because it was
    LLM-generated, nor is it less important to understand on the line-by-line
    level.
-   Anyone who reviews any line of code I commit can trust on my name and
    reputation that it was written with as much thought and concern as any other
    line.
-   Absolutely no prose (comments, documentation, PR bodies) authored by LLMs.
    The responsibility is on the human.

In a sense, I treat agentic LLMs as a glorified multi-file autocomplete. I am
not trusting it on any concepts it introduces, because 90% of the time I know
the concept already and probably know it better, or know of a better way to
approach or structure the problem. That's just how it is, empirically today, in
2026, for me personally.

My "fully vibed" code I strictly keep out of source control for repositories of
importance; they are kept in different repositories and appropriately version
controlled and recorded and noted with the circumstances in which they were
written. I have found the most mileage in these as simple test scripts ("write
an executable to test my library out on sample inputs", "grab the spectrogram
trace from the read-only API of my signal analyzer", "aggregate these results
into a quantitative notebook for me to reference") that are otherwise tedious to
write or involve dealing with throwaway APIs with low ROI in learning. In these
I also presume LLM comments are "load-bearing" for their effective
vibe-maintenance but what do I know.

Most of the thoughts in this post will be about the former category (projects of
importance), although sometimes my vibed scripts eventually transition into code
bases I treat with the same level of respect, where these strategies will also
apply.

## Haskell and LLMs

Now my personal opinion and wish: in an ideal world, Haskell *should* in agentic
software engineering as what LEAN is in agentic research mathematics: a
framework for LLMs to self-construct the scaffolding they need to guide
themselves to their correct goal.

Contrary to others, I do *not* believe that "correctness at generation-time" is
a plausible goal, not today in 2026, and probably not ever. The motion towards
correctness is asymptotic. You might get close enough for projects, but the long
tail of correctness is...long. And demands respect. That's why even the most
full-vibed frontier-model projects have [large suites of
tests](https://bun.com/blog/bun-in-rust) that exist to guide the agent. My bet
is that agentic coding in five years will not be "spit out the correct program
in one batch of text", it will be "set up the best scaffolding that guides the
implementation".

# Signoff

Hi, thanks for reading! You can reach me via email at <justin@jle.im>, or at
twitter at [\@mstk](https://twitter.com/mstk)! This post and all others are
published under the [CC-BY-NC-ND
3.0](https://creativecommons.org/licenses/by-nc-nd/3.0/) license. Corrections
and edits via pull request are welcome and encouraged at [the source
repository](https://github.com/mstksg/inCode).

If you feel inclined, or this post was particularly helpful for you, why not
consider [supporting me on Patreon](https://www.patreon.com/justinle/overview),
or a [BTC donation](bitcoin:3D7rmAYgbDnp4gp4rf22THsGt74fNucPDU)? :)

\end{document}

\end{document}
